\input{../tex-inputs/latex-leseansicht-vorspann}

\section[Arthur und Lili Schnitzler an Gustav Schwarzkopf, 29. 4. 1926]{L04476 Arthur und Lili Schnitzler an Gustav Schwarzkopf, 29. 4. 1926}
\nopagebreak\mylabel{L04476v}
\rehead{ }\normalsize\beginnumbering\briefempfaengerindex{Schwarzkopf, Gustav@\textsc{Schwarzkopf, Gustav}!zzzCappellini, Lili@\emph{von Lili Cappellini}!1926-04-292@{29. 4. 1926}|(be}\briefempfaengerindex{Schwarzkopf, Gustav@\textsc{Schwarzkopf, Gustav}!zzzSchnitzler, Arthur@\emph{von Arthur Schnitzler}!1926-04-292@{29. 4. 1926}|(be}
\toendnotes[C]{\smallbreak\pagebreak[2]}
\correspDesc{Versand  durch Arthur Schnitzler, Lili Schnitzler am 29. 4. 1926 in Las Palmas de Gran Canaria
\newline{}Übermittlung  am 30. 4. 1926 in Las Palmas de Gran Canaria
\newline{}Erhalt  durch Gustav Schwarzkopf im Zeitraum [30. 4. 1926 – 4. 5. 1926?] in Wien}\toendnotes[C]{\smallbreak}
\Standort{DLA, A:Schnitzler, HS.1985.1.1897.}
\physDesc{Bildpostkarte, 210 Zeichen
\newline{}Handschrift Arthur Schnitzler: Bleistift, lateinische Kurrent
\newline{}Handschrift Lili Cappellini: schwarze Tinte, lateinische Kurrent
\newline{}Versand: Stempel: »\nobreak{}\oindex{Las Palmas de Gran Canaria@\textbf{Las Palmas de Gran Canaria}|pwk}Las Palmas, 30-Abr-26, 11 M\nobreak{}«.  }\toendnotes[C]{\smallbreak}\pstart{}{\pb}\label{T_L04476-1v}\edtext{\textcolor{gray}{\textbf{A. S.}}}{\lemma{\textnormal{\emph{A. S.}}}\Cendnote{\textnormal{ovaler Absenderkleber}}}\label{T_L04476-1}\pend{}\pstart{}\textcolor{gray}{\textbf{WIEN, XVIII.}}\oindex{XVIII., Währing@\textbf{XVIII., Währing}|pw}\pend{}\pstart{}\textcolor{gray}{\textbf{STERNWARTESTR. 71}}\oindex{Wien@\textbf{Wien}!XVIII., Währing@\textbf{XVIII., Währing}!Sternwartestraße 71@\textbf{Sternwartestraße 71}|pw}\pend{}{\bigskip}\pstart{}Austria\oindex{Österreich@\textbf{Österreich}|pw}\pend{}\pstart{}Hrn Gustav Schwarzkopf\pend{}\pstart{}Vienne I\oindex{I., Innere Stadt@\textbf{I., Innere Stadt}|pw}\pend{}\pstart{}Tiefer Graben 17\oindex{Wien@\textbf{Wien}!I., Innere Stadt@\textbf{I., Innere Stadt}!Tiefer Graben 17@\textbf{Tiefer Graben 17}|pw}\pend{}{\bigskip}
\pstart
           \noindent{}\raggedleft{}{\pb}\textcolor{gray}{\textbf{Gran Canaria\oindex{Gran Canaria@\textbf{Gran Canaria}|pw} – Vega de Santa Brigida\oindex{Santa Brígida@\textbf{Santa Brígida}|pw}.}}\pend
           \vspace{1em}
\pstart
           {\pb}\label{T_L04476-2v}\edtext{San}{\lemma{\textnormal{\emph{San}}}\Cendnote{\textnormal{gut zu lesen, trotzdem vermutlich »Las« gemeint}}}\label{T_L04476-2} Palmas\oindex{Las Palmas de Gran Canaria@\textbf{Las Palmas de Gran Canaria}|pw} auf
               Gran Canaria\oindex{Gran Canaria@\textbf{Gran Canaria}|pw}, das also wirklich
      existirt. 29. 4. 26\pend
           \vspace{0.5em}
\pstart
           Herzliche Grüße Ihnen
               und den verehrten Brüdern\pwindex{Schwarzkopf, Emil 17.\,9.\,1851 Wien – 28.\,1.\,1928 ebd.@\textsc{Schwarzkopf, Emil} (17.\,9.\,1851 Wien – 28.\,1.\,1928 ebd.)|pwv}\pwindex{Schwarzkopf, Max 12.\,6.\,1857 Wien – 14.\,4.\,1928 ebd.@\textsc{Schwarzkopf, Max} (12.\,6.\,1857 Wien – 14.\,4.\,1928 ebd.)|pwv}.\pend
           
\pstart
           Märchenhaft schön!\pend
           \pstart Ihr \spacefill\mbox{Arthur}\pend{}\selectlanguage{ngerman}\vspace{1em}
\pstart
           \noindent{}{[}hs. Cappellini:{]} Schönste Grüße\pend
           \pstart \spacefill\mbox{Lili S.}\pend{}\selectlanguage{ngerman}\endnumbering\briefempfaengerindex{Schwarzkopf, Gustav@\textsc{Schwarzkopf, Gustav}!zzzCappellini, Lili@\emph{von Lili Cappellini}!1926-04-292@{29. 4. 1926}|)be}\briefempfaengerindex{Schwarzkopf, Gustav@\textsc{Schwarzkopf, Gustav}!zzzSchnitzler, Arthur@\emph{von Arthur Schnitzler}!1926-04-292@{29. 4. 1926}|)be}\mylabel{L04476h}
\begin{anhang}
\end{anhang}\newcommand{\dateiname}{L04476}\newcommand{\titel}{Arthur und Lili Schnitzler an Gustav Schwarzkopf, 29. 4. 1926}\newcommand{\editorInnen}{Herausgegeben von Jahnke, SelmaMüller, Martin Anton}\input{../tex-inputs/latex-leseansicht-abspann}
